\begin{problem}[Compute the stalk of the structure sheaf]\label{prob:vi17-p03}
For $\mfp\in\Spec A$, prove $\OO_{X,\mfp}\cong A_{\mfp}$.
\end{problem}

\ifdefined\FullProblemDossiers
\paragraph{Why this problem matters.}
The theorem identifies abstract sheaf germs with the local algebra already developed in commutative algebra.

\paragraph{Useful ideas.}
Principal opens containing $\mfp$ form a basis, so compute the stalk as a colimit of $A_f$ for $f\notin\mfp$.

\paragraph{Guided hints.}
\begin{enumerate}
\item Replace all neighborhoods by principal neighborhoods.
\item Every $a/s\in A_{\mfp}$ comes from the section $a/s\in A_s$.
\item Use the equality criterion for localization to prove injectivity.
\end{enumerate}

% VI17_FULL_SOLUTION_REFINED
% VI17_FULL_SOLUTION_REFINED
\begin{solution}
Principal opens $D(f)$ with $f\notin\mfp$ form a neighborhood basis of $\mfp$.  Using the principal-open computation, the stalk is therefore
\[
\OO_{X,\mfp}
\cong
\varinjlim_{f\notin\mfp}\OO_X(D(f))
\cong
\varinjlim_{f\notin\mfp}A_f.
\]
For every $f\notin\mfp$, the universal property of localization gives a map $A_f\to A_{\mfp}$; these maps are compatible and induce
\[
\Phi:\varinjlim_{f\notin\mfp}A_f\longrightarrow A_{\mfp}.
\]

The map is surjective because an arbitrary fraction $a/s\in A_{\mfp}$ has $s\notin\mfp$, hence already comes from the element $a/s\in A_s$.

For injectivity, suppose representatives $a/f^r\in A_f$ and $b/g^s\in A_g$ have the same image in $A_{\mfp}$.  By the equality criterion in $A_{\mfp}$ there exists $u\notin\mfp$ such that
\[
u(g^sa-f^rb)=0.
\]
In the common localization $A_{fgu}$ this equation says the two representatives are equal.  Thus they become equal at a later stage of the directed system and define the same colimit element.  Therefore $\Phi$ is injective, giving the canonical isomorphism
\[
\OO_{X,\mfp}\cong A_{\mfp}.
\]
\end{solution}



\paragraph{What happened mathematically.}
Taking a stalk means inverting every function that is nonzero at the chosen prime.

\paragraph{Extensions.}
Identify the maximal ideal and residue field of the stalk.

\paragraph{Related problems.}
VI.17.P2 and VI.17.P4 develop neighboring aspects of the same localization-sheaf calculus.

\paragraph{Provenance.}
\textbf{New synthesis of the stalk formula explicitly present in the legacy structure-sheaf discussion.}
\else
\paragraph{Provenance.} \textbf{New synthesis of the stalk formula explicitly present in the legacy structure-sheaf discussion.}
\fi
